% id: 81-12
% book: 베이직쎈 미적분1 · 05 도함수의 활용 (2)
% section: 유형 051 도함수의 그래프와 함수의 극대·극소
% figure: tikz:fig-81-12
% answer: 
% answer_source: 
% variant_level: 0 (원본 전사)
% 이 파일은 build.mjs 가 items.json 에서 만든다. 고칠 때는 items.json 을 고친다.
\smash{\raisebox{1.5mm}{\dmkichul{베이직쎈 미적분1 81쪽 12번}}}
\noindent\begin{minipage}[t]{0.58\linewidth}\vspace{0pt}\raggedright 함수 $f(x)$의 도함수 $y=f'(x)$의 그래프가 오른쪽 그림과 같을 때, 다음 중 옳지 \underline{않은} 것은?\end{minipage}\hfill\begin{minipage}[t]{0.4\linewidth}\vspace{0pt}\centering \resizebox{\ifdim\width>\linewidth\linewidth\else\width\fi}{!}{% 81-12(81쪽) — 도함수 $y=f'(x)$ 의 그래프($x=-3,\ -1,\ 2,\ 4$ 에서 $x$축과 만난다 · $x=-2,\ 3$ 에서 극소, $x=1$ 에서 극대). 스캔 화질이 낮아 TikZ 로 다시 그림.
% 원문에 있는 것만 옮긴다: 곡선 · $x=-2,\ 1,\ 3$ 의 안내 점선 · 눈금 -3, -2, -1(지시선), 1, 2, 3, 4 · 라벨 O, x, y, y=f'(x). 원문에 없는 함숫값은 적지 않는다.
\begin{tikzpicture}[x=0.36cm, y=0.32cm, line width=0.6pt]
  \draw[->, >=stealth, line width=0.7pt] (-4.2,0) -- (5.7,0) node[below=1pt, font=\small] {$x$};
  \draw[->, >=stealth, line width=0.7pt] (0,-2.75) -- (0,4.25) node[left=1pt, font=\small] {$y$};
  \draw[densely dotted, line width=0.4pt] (-2,0) -- (-2,-2.0);
  \draw[densely dotted, line width=0.4pt] (1,0) -- (1,2.45);
  \draw[densely dashed, line width=0.4pt] (3,0) -- (3,-2.25);
  \draw plot[smooth, tension=0.7] coordinates {(-3.75,3.4) (-3.5,2.3) (-3.2,0.9) (-3,0) (-2.7,-1.1) (-2.4,-1.8) (-2,-2.0) (-1.7,-1.6) (-1.4,-0.85) (-1,0) (-0.75,0.85) (-0.5,1.45) (-0.25,1.85) (0,2.08) (0.35,2.3) (0.7,2.43) (1,2.45) (1.25,2.3) (1.5,1.85) (1.75,1.0) (2,0) (2.25,-0.85) (2.55,-1.75) (2.8,-2.2) (3,-2.25) (3.3,-1.9) (3.6,-1.1) (3.85,-0.35) (4,0) (4.25,1.1) (4.5,2.4) (4.7,3.55)};
  \draw[->, >=stealth, line width=0.3pt] (-1.5,1.3) .. controls (-1.2,0.85) .. (-1.05,0.22);
  \node[below left=-1pt and 0pt, font=\small] at (-3,0) {$-3$};
  \node[above left=-2pt and -4pt, font=\small] at (-2,0) {$-2$};
  \node[anchor=south, font=\small] at (-1.85,1.35) {$-1$};
  \node[below left=-2pt and -3pt, font=\small] at (0,0) {$\pt{O}$};
  \node[below=1pt, font=\small] at (1,0) {$1$};
  \node[below=1pt, font=\small] at (2,0) {$2$};
  \node[above=1pt, font=\small] at (3,0) {$3$};
  \node[below=1pt, font=\small] at (4,0) {$4$};
  \node[anchor=south, font=\small] at (3.4,4.0) {$y=f'(x)$};
\end{tikzpicture}
}\end{minipage}\par
\par\medskip\par\noindent\hangindent=1.4em\hangafter=1 {\small ①}\ $f(x)$는 $x=-3$에서 극댓값을 갖는다.\par\noindent\hangindent=1.4em\hangafter=1 {\small ②}\ $f(x)$는 구간 $(-2{,}\mkern3mu\mathopen{}-1)$에서 감소한다.\par\noindent\hangindent=1.4em\hangafter=1 {\small ③}\ $f(x)$는 구간 $(1{,}\mkern3mu\mathopen{}2)$에서 증가한다.\par\noindent\hangindent=1.4em\hangafter=1 {\small ④}\ $f(x)$는 $x=4$에서 극솟값을 갖는다.\par\noindent\hangindent=1.4em\hangafter=1 {\small ⑤}\ 구간 $(-3{,}\mkern3mu\mathopen{}4)$에서 $f(x)$가 극값을 갖는 $x$의 값은 3개이다.\par\medskip
